\indent ಇಲ್ಲಿ ಕೊಟ್ಟಿರುವ ವಾಕ್ಯಗಳಲ್ಲಿ ಕಂಡುಬರುವ ಕ್ರಿಯಾಪದಗಳನ್ನು ಗಮನಿಸಿರಿ.\par

\indent (ಪ್ರ. ಪು. ಏ.) ... ಅವನು ಓಡುತ್ತಾನೆ - ಓಡು + ಉತ್ತ + ಆನೆ\par

\indent (ಪ್ರ. ಸ್ತ್ರೀ. ಏ.) ... ಅವಳು ಓಡುತ್ತಾಳೆ - ಓಡು + ಉತ್ತ + ಆಳೆ\par

\indent (ಪ್ರ. ನ. ಏ.) ... ಅದು ಓಡುತ್ತದೆ - ಓಡು + ಉತ್ತ + ಅದೆ\par

\indent (ಮ. ವಾ. ಏ.) ... ನೀನು ಓಡುತ್ತೀಯೆ - ಓಡು + ಉತ್ತ + ಈಯೆ\par

\indent (ಉ. ವಾ. ಏ.) ... ನಾನು ಓಡುತ್ತೇನೆ - ಓಡು + ಉತ್ತ + ಏನೆ\par

\indent (ಪ್ರ. ವಾ. ಬ.) ... ಅವರು ಓಡುತ್ತಾರೆ - ಓಡು + ಉತ್ತ + ಆರೆ\par

\indent (ಪ್ರ. ನ. ಬ.) ... ಅವು ಓಡುತ್ತವೆ - ಓಡು + ಉತ್ತ + ಅವೆ\par

\indent (ಮ. ವಾ. ಬ.) ... ನೀವು ಓಡುತ್ತೀರಿ - ಓಡು + ಉತ್ತ + ಈರಿ\par

\indent (ಉ. ವಾ. ಬ.) ... ನಾವು ಓಡುತ್ತೇವೆ - ಓಡು + ಉತ್ತ + ಏವೆ\par

\indent ಇಲ್ಲಿ ಕ್ರಿಯಾಪದಗಳಿವೆ. ಈ ಕ್ರಿಯಾಪದಗಳಲ್ಲಿ “ಓಡು” ಧಾತುವಿನ ಮುಂದೆ ಎಲ್ಲ ಕಡೆಗೂ “ಉತ್ತ” ಎಂಬ ಪ್ರತ್ಯಯ ಬಂದಿದೆ ; ಈ ಪ್ರತ್ಯಯದ ಮುಂದೆ ಏಕವಚನದಲ್ಲಿ ಆನೆ, ಆಳೆ, ಅದೆ, ಈಯೆ, ಏನೆ ಎಂಬ ಆಖ್ಯಾತ ಪ್ರತ್ಯಯಗಳೂ ಬಹುವಚನದಲ್ಲಿ ಆರೆ, ಅವೆ, ಈರಿ, ಏವೆ ಎಂಬ ಆಖ್ಯಾತ ಪ್ರತ್ಯಯಗಳೂ ಬಂದಿವೆ.\par

\indent ಮೂಲ ಧಾತುವಿಗೂ ಆಖ್ಯಾತ ಪ್ರತ್ಯಯಕ್ಕೂ ನಡುವೆ ಬಂದಿರುವ “ಉತ್ತ” ಎಂಬುದು ಕಾಲ ಸೂಚಕ ಪ್ರತ್ಯಯ. ಇದು ವರ್ತಮಾನ ಕಾಲ ವನ್ನು ಸೂಚಿಸುವುದು. ಅಂದರೆ, ಕ್ರಿಯಾಪದವು ವರ್ತಮಾನ ಕಾಲದಲ್ಲಿರುವು ದೆಂಬುದನ್ನು ತಿಳಿಸುವುದು.\par

\kbvsechead{ಸಂಕೇತಾಕ್ಷರಗಳ ನಿರ್ಧಾರ :}

\begin{kbvderiv}
\vspace{4pt plus 3pt minus 1pt}
\noindent{\renewcommand{\arraystretch}{1.24}\setlength{\tabcolsep}{5pt}
\begin{tabularx}{\textwidth}{X X X}
ಪ್ರ. - ಪ್ರಥಮ ಪುರುಷ & ಪು - ಪುರುಷ & ಏ. - ಏಕವಚನ \\
ಮ. - ಮಧ್ಯಮ ಪುರುಷ & ಸ್ತ್ರೀ. - ಸ್ತ್ರೀಲಿಂಗ & ಬ. - ಬಹುವಚನ \\
ಉ. - ಉತ್ತಮ ಪುರುಷ & ವಾ. - ವಾಚ್ಯಲಿಂಗ &  \\
\end{tabularx}}
\vspace{4pt plus 3pt minus 1pt}
\end{kbvderiv}
